\subsection{Supply availability and the storage contribution}\label{sec:supply-results}

The additional supply experiment evaluates previously frozen gains, with a newly declared boundary-resource measurement. It does not replace the four-objective negative comparison in Section 4.9. Table~\ref{tab:supply-normal} and Figure~\ref{fig:supply-normal} report all four normal-supply seeds; the availability hybrid is numerically identical to the scalar hybrid on these paths.

\begin{table}[H]\centering\small
\caption{New normal-supply repeated-load paths, seeds 5500--5503, 1 ms RK4. Means use the common 8.01--1800 s window. B is the declared boundary-resource index, not total input energy. All four paths enter holding under each policy.}\label{tab:supply-normal}
\begin{tabularx}{\linewidth}{@{}>{\raggedright\arraybackslash}p{.24\linewidth}*{5}{>{\centering\arraybackslash}X}@{}}\toprule
Policy & Pass & RMSE (mm) & $V_s$ (mL) & Oil (mL) & $B$ (kJ)\\\midrule
Scalar hybrid & 4/4 & 1.271 & 1085.0 & 21.55 & 23.08\\
Active PID & 4/4 & 0.314 & 1592.7 & 34.58 & 33.93\\
Active SMC & 4/4 & 0.252 & 13405.7 & 13.49 & 281.71\\
Availability hybrid & 4/4 & 1.271 & 1085.0 & 21.55 & 23.08\\
\bottomrule\end{tabularx}\end{table}

Scalar hybrid/active SMC mean main-supply volume is 1085.0/13405.7 mL and mean B is 23.08/281.71 kJ. Their RMSE is 1.271/0.252 mm. The relative main-supply-volume difference, 91.91\%, is a ratio of these group means. It quantifies a resource--tracking trade-off at frozen gains; it is not an energy-saving rate or evidence that SMC cannot reach a different trade-off after resource-oriented retuning. S12 gives the paired differences. The earlier four-objective ordering favoring SMC remains valid.

\begin{figure}[H]\centering\includegraphics[width=\linewidth]{risk_extension_supply_normal.pdf}\caption{Normal-supply resource and tracking comparison. Small markers show all four seeded paths; large markers show means. The two hybrid instances coincide because the availability extension is inactive. B includes main-source and storage-outlet work only; no pump-input or efficiency claim follows.}\label{fig:supply-normal}\end{figure}

For the eight new joint vectors with supply disconnected on [300,1500) s, complete-path pass counts are Scalar hybrid: 8/8 with storage and 1/8 without; Active PID: 8/8 with storage and 1/8 without; Active SMC: 8/8 with storage and 0/8 without; Availability hybrid: 8/8 with storage and 1/8 without. Figure~\ref{fig:supply-interruption} includes every on/off outcome; passing requires hold entry, both 10 mm screens, and no relief or stroke event. The current availability signal is ideal and all policies retain the same return path.

Within this storage-enabled joint family, scalar hybrid/active SMC mean main-source volume is 539.9/4903.6 mL and B is 15.02/110.12 kJ. These separately paired resource comparisons are reported in S12; their schedules differ from the normal repeated-load family.

The availability hybrid's worst sampled post-transition error is 9.215 mm and worst confirmed-hold drop is 2.051 mm. Mean within-interruption maximum error is 1.173 mm, compared with 3.571 mm for the scalar hybrid; the complete-path screen still includes the earlier transition. Mean storage use is 409.74/286.68 mL, respectively. Paired changes are given in S12. This small fixed design evaluates the stated availability rule and does not establish a superior controller family.

At fixed plant and noise inputs, enabling storage changes a failed complete path into a passing path in Scalar hybrid: 7/8, Active PID: 7/8, Active SMC: 8/8, Availability hybrid: 7/8. The reverse changes, unchanged pass/fail cases and exact seeds are also retained in S12. Storage necessity is therefore judged for each declared policy and interruption, while the original continuously supplied disabled-storage successes remain valid. A common pass ratio cannot establish physical reliability or universal necessity.

\begin{figure}[H]\centering\includegraphics[width=\linewidth]{risk_extension_supply_interruption.pdf}\caption{All eight paired joint vectors during the 1800 s supply-availability study. Each value is the larger of maximum post-transition absolute error and confirmed-hold drop; both have the same 10 mm limit. Thus an error below 10 mm alone cannot hide a drop failure. Maximum error includes the pre-interruption transition. Other pass conditions and all failed seeds are reported in S12.}\label{fig:supply-interruption}\end{figure}

The prespecified seed-5400 refinement gives a largest absolute drop change of 0.212 mm and maximum-error change of 0.195 mm across the eight policy/storage pairs. Pass status changes in 0 of eight pairs. This is a representative numerical sensitivity check, not convergence certification over all joint vectors.

The archived both-port-isolation exploration provides an adverse mechanism check: all eight 600 s interruption paths fail the 10 mm screen, including storage-enabled paths. Added storage alone cannot supply the required return balance, and gas--A pressure equalization can stop useful discharge while oil remains. That failed boundary is distinct from Equation~\eqref{eq:supply-flow}. The prospectively frozen supply-only results do not repair or replace those failures.
